\section{Bridge to recovery and sequential control}
\label{sec:b-bridge}

\paragraph{Recovery (optional extension).}
A recovery action $g_r$ becomes a first-study candidate only with a stable
precondition, a declared rollback boundary, a repair action and a continuation
contract.  It must then demonstrate an increment over continue/verify under the
same branch budget and information, and it is evaluated as a third action with
its own feasibility mask and cost vector.  A generic unconstrained ``try again''
is not comparable and is not reported as recovery.

\paragraph{Sequential control and DAC/RL (future work).}
Extending the decision to every checkpoint changes the problem in three ways
that must be earned rather than assumed.  First, the continuation is no longer
frozen: the value of acting at $t$ depends on the policy that acts later, and
the branch design no longer supplies a single backup.  Second, coverage becomes
policy-dependent: pessimistic value iteration requires uniform confidence
validity over the states the new policy reaches, and those include
post-intervention states that depth-one branching never re-enters; the one-step
result of \cref{prop:transfer,rem:transfer-fails} makes this gap explicit and
shows a registered-scale construction in which the repeated controller loses
while every one-step quantity is positive.  Third, the allocation problem
becomes a sequential decision under a budget, i.e.\ a constrained
semi-Markov decision process over checkpoints with option durations and cost
accumulation; the one-step problem in this paper is its first backup, and the
constraint structure is a natural fit for constrained RL or decision-focused
allocation methods.  None of that is claimed here.  A depth-two branching pilot
that measures the post-intervention span would be the first instrument
requirement of that program, and it is not part of this paper.
